% Methods: platform design, the five methods, the gate, and the
% autonomous per-commit generation methodology. Instruction only.

[PROCESSING INSTRUCTION (Methods). Describe, as a senior author would,
how the platform is built and how it was generated. Cover four things in
order, each as a short subsection.

(1) The five established methods, one daily deliverable. Explain the four
producing stages and the orchestrator, then the fifth method, the VVUQ
gate, that clears each candidate. Ground every claim in pipeline/ source.

(2) The VVUQ gate and its thresholds. State the three dimensions and the
exact default thresholds: verification fraction must equal 1.0 (six
checks, schema valid, lint clean, 200K file cap); validation agreement at
least 0.95 with maximum relative error at most 0.05 and a recorded human
review; uncertainty at least 3 runs with coefficient of variation at most
0.10. Note the gate blocks on any failure and escalates divergence.
Ground in vvuq/ and configs/vvuq\_thresholds.yaml.

(3) Stage 1 automation and the Stage 2 references. Briefly: triple
simulation with consensus, robust web and PDF ingestion, 200K autochunk
with reconstruction READMEs, and the non-stop commit scheduler; then the
Stage 2 lights-off factory and hybrid surgery and medicine pilot, disabled
by default and gated on human oversight and a reachable VVUQ gate.

(4) Autonomous per-commit generation. Explain that Claude Code Opus 4.7
(1M context) Max generated the codebase in a managed cloud container, one
section per commit, pushed in real time within a single pull request, then
re-ran every CI check and fixed issues before a final repository-update
commit. Present the codegen commit metrics in
Table~\ref{tab:methods_codegen_commits} as evidence of scalable, real
life, per commit LLM VVUQ processing, and contrast with prior non
incremental systems.

SOURCE FILES:

\begin{table}[H]
\caption{Methods source map (abbreviated paths under repo root).}
\label{tab:methods_sources}
\centering
\begin{tabular}{L{4.6cm} L{8.8cm}}
\toprule
\textbf{Source (under repo root)} & \textbf{Use for} \\
\midrule
  pipeline/ : 6 stage modules + examples-pipeline/ & The four producing stages and the orchestrator; the daily deliverable model. \\
  vvuq/ : 4 assurance modules + examples-vvuq/ & The fifth method, the gate; the three VVUQ dimensions. \\
  configs/ : pipeline + vvuq YAML configs & The 200K cap, the 3 run default, and the gate thresholds. \\
  physical-ai/ : factory + pilot modules, examples & Stage 2 references and the safety interlocks. \\
  simulation/, ingestion/, chunking/, scheduler/ & Stage 1 automation around the pipeline. \\
  scripts/ verify script; tests/ : 8 modules, conftest & The environment gate and the verification floor. \\
  execution/01-foundation/ : env, tests, lint reports & Methods reproducibility: Python 3.11.15, tooling, clean static analysis. \\
\bottomrule
\end{tabular}
\end{table}

SOURCE DATA to render as a publication table (do not alter the numbers):

\begin{table}[H]
\caption{Autonomous codegen commits that built the platform.}
\label{tab:methods_codegen_commits}
\centering
\begin{tabular}{C{0.7cm} L{2.0cm} L{5.4cm} C{1.6cm} C{1.6cm}}
\toprule
\textbf{\#} & \textbf{SHA} & \textbf{Section} & \textbf{Gap} & \textbf{Churn} \\
\midrule
  C1 & c9e1404 & Foundation: tooling, CI, community & - & +1088 -2 \\
  C2 & a111ca8 & Pipeline (5 established methods) & 4m42s & +825 \\
  C3 & 187ba9e & VVUQ framework & 2m44s & +691 \\
  C4 & dcb11c0 & Stage 1 automation (4 modules) & 4m46s & +1127 \\
  C5 & 082fa9f & Stage 2 deployment refs & 1m49s & +389 \\
  C6 & 13b253d & Fix all errors (2nd to last) & 0m48s & +3 -1 \\
  C7 & bb27fd6 & Repo updates (last) & 3m05s & +119 -9 \\
\bottomrule
\end{tabular}
\end{table}

Also render the supporting breakdowns as compact tables: files by type
(44 .py / 2829 lines; 24 .md / 1140; 2 .yaml / 79; 1 .yml / 73; 1 .txt /
43; 1 .toml / 32; 1 .cff / 30; LICENSE / 21; .gitignore / 6; 76 files
total including the pre-existing LICENSE) and directory lines (pipeline
753, vvuq 573, tests 559 over 53 tests, physical-ai 325, ingestion 278,
simulation 266, chunking 260, .github 185, scheduler 143, scripts 112,
configs 79). Reference these tables; do not narrate every row.

HOW TO PROCESS: explain the architecture in prose and let the tables carry
the file inventory; abbreviate long module names. Reuse the
configs/vvuq\_thresholds.yaml numbers verbatim. Tie the per commit cadence
to the thesis by noting that the assurance, not the generation, dominated
the effort. Cite \cite{repo-cancer-automated}, \cite{anthropicopus47},
\cite{anthropicclaudecodeweb}, \cite{iec62304}, and \cite{iche6r3}.]
